\section{Existing Solution Introduction}
Recent advancements in artificial intelligence have led to the development of powerful Multimodal Large Language Models (MLLMs) such as GPT-4V and Gemini Pro Vision, which demonstrate remarkable capabilities in understanding visual content. However, these models are primarily optimized for short-form video snippets or static images. When applied to long-form video intelligence, they suffer from extreme computational overhead and a lack of temporal continuity, limiting their utility for complex, multi-hour narrative analysis.

On the other hand, traditional video analytics platforms offer functionalities like object detection, facial recognition, and motion tracking. Yet, they lack natural language conversational interfaces, multimodal semantic fusion (vision + audio + OCR), and the ability to synthesize high-level narratives from raw detection logs. Very few existing tools combine spatial detection, temporal entity tracking, and recursive summarization into a single unified framework. This technological gap has prompted increasing interest in \textbf{Retrieval-Augmented Generation (RAG)} for video, which combines dense multimodal retrieval with generative models to ground responses in verifiable sensor data.

\section{Relevant Theories and Models}
The foundation of modern video intelligence assistants lies in the fusion of Computer Vision (CV), Automatic Speech Recognition (ASR), and Large Language Models (LLMs). 

\subsection{Vision and Spatial Extraction}
Retrieval-Augmented Generation (RAG) for video relies on robust spatial feature extractors. The \textbf{YOLO (You Only Look Once)} framework \cite{yolov8}, specifically YOLOv8, frames object detection as a regression problem, allowing for real-time identification of entities with high Mean Average Precision (mAP). Complementary to this is \textbf{CLIP (Contrastive Language-Image Pre-training)} \cite{clip}, which maps visual and textual information into a unified latent space, enabling semantic scene-anchoring without the need for fixed-class labels.

\subsection{Audio Intelligence and ASR}
Automatic Speech Recognition has been revolutionized by Transformer-based models like \textbf{OpenAI's Whisper} \cite{whisper}. Trained on 680,000 hours of multilingual audio, Whisper provides high-fidelity timestamped transcriptions that are essential for grounding video narratives in spoken dialogue. Its robustness to background noise makes it a primary choice for local-first execution.

\subsection{GraphRAG and Knowledge Graphs}
The emergence of \textbf{GraphRAG} \cite{graphrag} represents a shift from flat vector stores to relational knowledge bases. By representing data as a graph $G = \langle V, E \rangle$, where $V$ denotes entities (people, objects, OCR text) and $E$ denotes temporal interactions, systems can perform "multi-hop" reasoning. This is critical for solving the problem of entity persistence across disjoint video segments.

\subsection{Recursive Map-Reduce Logic}
Hierarchical summarization is supported by \textbf{Map-Reduce} strategies \cite{mapreduce}. In this paradigm, long-form data is "mapped" into dense local summaries and subsequently "reduced" through recursive consolidation. This prevents the "Context Window Wall" encountered when feeding bulk metadata directly into an LLM.

\subsection{Vector Similarity Metrics and Retrieval Logic}
The efficiency of RAG frameworks is predicated on the accuracy of the retrieval phase. Standard methodologies utilize \textbf{Cosine Similarity} to measure the angular distance between a query vector $q$ and a document vector $d$ in a high-dimensional embedding space. Mathematically, this is expressed as:
\[ \text{similarity} = \cos(\theta) = \frac{\mathbf{A} \cdot \mathbf{B}}{\|\mathbf{A}\| \|\mathbf{B}\|} \]
While effective for broad semantic matching, Cosine Similarity often fails to capture the fine-grained temporal relationships required for video forensics. VidChain addresses this by integrating a multi-stage retrieval pipeline.

\subsection{Cross-Encoder Reranking Algorithms}
Research indicates that \textbf{Bi-Encoders} (used by vector stores like ChromaDB) provide high-speed search but suffer from lower relational accuracy. To bridge this gap, modern architectures implement \textbf{Cross-Encoders} for the final reranking phase. Unlike Bi-Encoders, which process query and document independently, Cross-Encoders (e.g., \textbf{MS-MARCO}) feed the concatenated query-document pair into a transformer simultaneously. This allows the model to capture deep interactions between the user's intent and the retrieved video segment, significantly reducing false-positive retrieval.

Assessing the quality of generated video narratives requires standardized quantitative benchmarks. The literature primarily identifies three metrics: \textbf{BLEU (Bilingual Evaluation Understudy)}, which measures $n$-gram overlap; \textbf{ROUGE (Recall-Oriented Understudy for Gisting Evaluation)}, which focuses on summary coverage; and \textbf{METEOR}, which addresses synonyms and morphological variants to provide a more human-aligned accuracy score. Together, these metrics ensure that narrative synthesis is both factually grounded and linguistically coherent.

\section{Comparison of Similar Projects}
Several systems have attempted to bridge the gap between video data and natural language reasoning, most notably end-to-end models like \textbf{Video-ChatGPT} and \textbf{Video-LLaVA}. While these models excel at short-form visual Q\&A on benchmarks such as \textbf{ActivityNet-QA}, they face significant hurdles in professional forensic contexts. Primarily, their \textbf{High Resource Consumption} necessitates server-grade GPUs (40GB--80GB VRAM), effectively prohibiting edge-local deployment. Furthermore, without a structured tracking layer, these monolithic systems are prone to \textbf{Temporal Hallucination} and the "Lost in the Middle" phenomena, where the model loses focus on critical chronological details as video length increases beyond its fixed context window.

\subsection{VidChain: A Hybrid Paradigm}
In contrast to the monolithic approach, VidChain focuses on a \textbf{Hybrid Logic}. It combines the high-speed retrieval of ChromaDB with the structural persistence of **GraphRAG**. This richer framework allows for more transparent, explainable, and robust academic assistance—particularly useful for longitudinal studies and technical explanations where traceability is paramount.

\section{Identified Gaps in the Literature}
Despite the rapid evolution of multimodal RAG, several critical gaps persist in the current research landscape. Most existing frameworks lack \textbf{Edge-Optimized Orchestration}, creating a problematic dependency on high-latency cloud APIs that prevents local execution on consumer-grade hardware. Furthermore, there is a notable absence of \textbf{Temporal Gating} mechanisms capable of managing computational loads during the fusion process. From a reasoning perspective, the \textbf{Entity Persistence Gap} remains a major hurdle; standard Video RAG systems often lack the structural memory required for multi-hop relational reasoning, such as tracking a specific subject's reappearances across temporal segments. Finally, a pervasive lack of \textbf{Citation Transparency} in many systems results in black-box summaries that fail to provide verifiable links to the specific visual or audio evidence justifying their claims. VidChain is specifically designed to bridge these gaps through its hybrid vector-graph architecture.

\section{Project Timeline and Methodology}
\label{sec:timeline}

The development of VidChain followed a structured sequential methodology to ensure the integrity of the forensic pipeline. The project was divided into eight distinct phases, spanning from initial planning to final submission. The following Gantt chart (Figure~\ref{fig:gantt}) illustrates the temporal distribution of these efforts across the Spring 2026 semester, reflecting the high computational complexity of the implementation and testing phases.

\begin{figure}[H]
\centering
\begin{ganttchart}[
    hgrid,
    vgrid={*1{black!10}},
    x unit=0.6cm,
    y unit chart=0.6cm,
    title height=1,
    title label font=\tiny\bfseries,
    bar label font=\footnotesize,
    bar/.style={fill=blue!60},
    bar height=0.5,
    progress label text={},
    canvas/.style={draw=none}
]{1}{12}
    \gantttitle{Spring 2026 Semester Timeline}{12} \\
    \gantttitle{Feb 02}{2} \gantttitle{Feb 16}{2} \gantttitle{Mar 02}{2} \gantttitle{Mar 16}{2} \gantttitle{Mar 30}{2} \gantttitle{Apr 13}{2} \\
    \gantttitle{2026}{2} \gantttitle{2026}{2} \gantttitle{2026}{2} \gantttitle{2026}{2} \gantttitle{2026}{2} \gantttitle{2026}{2} \\
    \ganttbar{1. Project Planning}{1}{1} \\
    \ganttbar{2. Requirement Engineering}{2}{2} \\
    \ganttbar{3. Data Collection \& Exploration}{3}{4} \\
    \ganttbar{4. System Design}{5}{5} \\
    \ganttbar{5. Implementation (Nodes \& Pipeline)}{6}{10} \\
    \ganttbar{6. Testing, Debugging \& Evaluation}{11}{11} \\
    \ganttbar{7. Documentation}{11}{12} \\
    \ganttbar{8. Final Review \& Submission}{12}{12}
\end{ganttchart}
\caption{Project Development Timeline (Finalized April 26, 2026)}
\label{fig:gantt}
\end{figure}

\section{Conclusion of Survey}
The literature review highlights a dynamic landscape of video intelligence tools. While general-purpose models offer linguistic fluency, they lack the scholarly grounding and resource efficiency required for edge-local video analysis. This survey reinforces the importance of a decentralized research assistant that integrates real-time retrieval, semantic reranking, and citation transparency. The identified limitations in existing solutions provide a strong justification for the proposed **VidChain** system, which seeks to bridge these gaps and enhance the reliability of video-based information discovery.
